\documentclass[11pt,a4paper]{article}

\def\courseTitle{Industrial Security (Master)}
\def\labNumber{07}
\def\labTitle{Safety / Security Co-Engineering on a SIL-rated SIS}
\def\labSection{07}
\def\labTime{4 hours, group of 4--6}

% =====================================================================
% Shared preamble for lab sheets.
%
% Caller must \documentclass{article} first, then define:
%   \def\courseTitle{...}
%   \def\labNumber{NN}
%   \def\labTitle{...}
%   \def\labTime{e.g. "30--45 min"}
%   \def\labSection{e.g. "02"}
% Then % =====================================================================
% Shared preamble for lab sheets.
%
% Caller must \documentclass{article} first, then define:
%   \def\courseTitle{...}
%   \def\labNumber{NN}
%   \def\labTitle{...}
%   \def\labTime{e.g. "30--45 min"}
%   \def\labSection{e.g. "02"}
% Then % =====================================================================
% Shared preamble for lab sheets.
%
% Caller must \documentclass{article} first, then define:
%   \def\courseTitle{...}
%   \def\labNumber{NN}
%   \def\labTitle{...}
%   \def\labTime{e.g. "30--45 min"}
%   \def\labSection{e.g. "02"}
% Then \input{../../template/lab-preamble.tex}
% =====================================================================

\usepackage[utf8]{inputenc}
\usepackage[T1]{fontenc}
\usepackage[english]{babel}
\usepackage[a4paper, margin=2cm]{geometry}
\usepackage{graphicx}
\usepackage{listings}
\usepackage{xcolor}
\usepackage{colortbl}
\usepackage{tikz}
\usepackage{amsmath}
\usepackage{amssymb}
\usepackage{booktabs}
\usepackage{enumitem}
\usepackage{titlesec}
\usepackage{fancyhdr}
\usepackage{hyperref}
\usepackage{tcolorbox}
\tcbuselibrary{skins, breakable}
\usepackage{fontawesome5}

\input{../../../template/tikz-styles.tex}

\providecommand{\courseAuthor}{Industrial Security Lecture Series}
\providecommand{\courseInstitute}{Department of Computer Science}
\providecommand{\companyLogo}{logo}

\definecolor{slideAccent}{HTML}{00205B}
\definecolor{slideSecondary}{HTML}{00A0DC}

% --- Corporate branding overrides ------------------------------------
\IfFileExists{../../../branding.tex}{\input{../../../branding.tex}}{}

\colorlet{labAccent}{slideAccent}
\colorlet{labSec}{slideSecondary}
\definecolor{labCheck}{HTML}{2E7D32}
\definecolor{labWarn}{HTML}{B23A48}

\lstset{
  backgroundcolor=\color{black!4},
  basicstyle=\ttfamily\footnotesize,
  breaklines=true,
  frame=leftline,
  framerule=2pt,
  rulecolor=\color{labAccent},
  numbers=left,
  numberstyle=\tiny\color{black!50},
  showstringspaces=false,
  tabsize=2
}

\setlength{\tabcolsep}{4pt}

% Let TeX add a little extra inter-word stretch as a last resort before
% it reports an Overfull \hbox. Only kicks in on lines that would
% otherwise overflow (long \texttt paths, URLs, CIDRs); never loosens a
% line that already fits. Keeps prose/itemize within the margin.
\setlength{\emergencystretch}{3em}

\pagestyle{fancy}
\fancyhf{}
\renewcommand{\headrulewidth}{0.4pt}
\lhead{\small \courseTitle{} -- Lab \labNumber}
\rhead{\small Maps to Section \labSection}
\cfoot{\thepage}

\newtcolorbox{stepbox}[1]{
  enhanced, breakable, arc=2pt, boxrule=0pt,
  colback=labAccent!7, colframe=labAccent, leftrule=3pt,
  fonttitle=\bfseries\small, coltitle=white,
  colbacktitle=labAccent,
  title={\faTools\ Step #1}
}

\newtcolorbox{warnbox}{
  enhanced, breakable, arc=2pt, boxrule=0pt,
  colback=labWarn!7, colframe=labWarn, leftrule=4pt,
  fonttitle=\bfseries\small, coltitle=white, colbacktitle=labWarn,
  title={\faExclamationTriangle\ Caution}
}

\newtcolorbox{notebox}{
  enhanced, breakable, arc=2pt, boxrule=0pt,
  colback=labAccent!4, colframe=labAccent, leftrule=2pt,
  fonttitle=\bfseries\small, coltitle=white, colbacktitle=labAccent,
  title={\faInfoCircle\ Note}
}

\newtcolorbox{goalbox}{
  enhanced, breakable, arc=2pt, boxrule=0pt,
  colback=labCheck!7, colframe=labCheck, leftrule=3pt,
  fonttitle=\bfseries\small, coltitle=white, colbacktitle=labCheck,
  title={\faBullseye\ Goal}
}

\newtcolorbox{deliverbox}{
  enhanced, breakable, arc=2pt, boxrule=0pt,
  colback=labSec!10, colframe=labSec, leftrule=3pt,
  fonttitle=\bfseries\small, coltitle=white, colbacktitle=labSec,
  title={\faClipboardList\ Deliverable}
}

% --- Solution-sheet boxes (used only by lab-solutions.tex) ----------
\newtcolorbox{solutionbox}[1]{
  enhanced, breakable, arc=2pt, boxrule=0pt,
  colback=labCheck!7, colframe=labCheck, leftrule=3pt,
  fonttitle=\bfseries\small, coltitle=white, colbacktitle=labCheck,
  title={\faCheckCircle\ #1}
}

\newtcolorbox{rubricbox}{
  enhanced, breakable, arc=2pt, boxrule=0pt,
  colback=labSec!7, colframe=labSec, leftrule=3pt,
  fonttitle=\bfseries\small, coltitle=white, colbacktitle=labSec,
  title={\faBalanceScale\ Grading rubric}
}

\titleformat{\section}{\large\bfseries\color{labAccent}}{\thesection}{0.5em}{}

\newcommand{\labheader}{%
  \noindent
  \begin{tikzpicture}
    \fill[labAccent] (0,0) rectangle (\textwidth, -1.6cm);
    \node[anchor=north west, text=white, font=\bfseries\large] at (0.6cm, -0.4cm) {Lab \labNumber{} -- \labTitle};
    \node[anchor=south west, text=white, font=\small] at (0.6cm, -1.3cm) {\courseTitle{} \quad $\bullet$ \quad Maps to Section \labSection{} \quad $\bullet$ \quad \labTime};
    \node[anchor=north east, fill=labSec, text=white, font=\bfseries, inner sep=4pt, rounded corners=2pt]
      at ([xshift=-0.4cm, yshift=-0.4cm]\textwidth,0) {LAB \labNumber};
  \end{tikzpicture}
  \vspace{4mm}
}

% =====================================================================

\usepackage[utf8]{inputenc}
\usepackage[T1]{fontenc}
\usepackage[english]{babel}
\usepackage[a4paper, margin=2cm]{geometry}
\usepackage{graphicx}
\usepackage{listings}
\usepackage{xcolor}
\usepackage{colortbl}
\usepackage{tikz}
\usepackage{amsmath}
\usepackage{amssymb}
\usepackage{booktabs}
\usepackage{enumitem}
\usepackage{titlesec}
\usepackage{fancyhdr}
\usepackage{hyperref}
\usepackage{tcolorbox}
\tcbuselibrary{skins, breakable}
\usepackage{fontawesome5}

% =====================================================================
% Shared TikZ styles for the Industrial Security lecture series.
% Loaded by every slide deck and exercise sheet that uses TikZ.
% =====================================================================

\usetikzlibrary{
  arrows.meta,
  shapes,
  shapes.geometric,
  shapes.symbols,
  shapes.multipart,
  positioning,
  fit,
  calc,
  backgrounds,
  decorations.pathmorphing,
  decorations.pathreplacing,
  decorations.text,
  patterns,
  matrix,
  shadows
}

% --- colour palette --------------------------------------------------
\definecolor{isOT}{HTML}{1F4E79}     % OT (deep blue)
\definecolor{isIT}{HTML}{6F2DA8}     % IT (purple)
\definecolor{isSafety}{HTML}{C00000} % safety red
\definecolor{isNet}{HTML}{2E7D32}    % network green
\definecolor{isAttack}{HTML}{B23A48} % attack red
\definecolor{isAccent}{HTML}{E08E0B} % amber
\definecolor{isMute}{HTML}{6B6B6B}   % grey
\definecolor{isLight}{HTML}{EDF2F8}  % light background
\definecolor{isPLC}{HTML}{2C5282}
\definecolor{isHMI}{HTML}{2F855A}
\definecolor{isSCADA}{HTML}{6B46C1}

% --- generic node styles --------------------------------------------
\tikzset{
  isnode/.style={
    draw, rounded corners=2pt, minimum height=8mm, minimum width=18mm,
    font=\small, align=center, thick
  },
  isbox/.style={isnode, fill=isLight},
  plc/.style={isnode, fill=isPLC!15, draw=isPLC, thick},
  hmi/.style={isnode, fill=isHMI!15, draw=isHMI, thick},
  scada/.style={isnode, fill=isSCADA!15, draw=isSCADA, thick},
  sensor/.style={isnode, fill=isNet!10, draw=isNet, minimum height=6mm, minimum width=14mm, font=\scriptsize},
  actuator/.style={isnode, fill=isAccent!15, draw=isAccent, minimum height=6mm, minimum width=14mm, font=\scriptsize},
  firewall/.style={isnode, fill=isAttack!15, draw=isAttack, thick},
  server/.style={isnode, fill=isMute!15, draw=isMute!80, thick},
  switch/.style={isnode, fill=isLight, draw=black!60, minimum height=6mm, minimum width=14mm, font=\scriptsize},
  workstation/.style={isnode, fill=white, draw=isOT, thick},
  attacker/.style={isnode, fill=isAttack!20, draw=isAttack, thick, font=\small\bfseries},
  inetnode/.style={draw, ellipse, minimum width=24mm, minimum height=10mm, font=\scriptsize, fill=isLight, thick},
  process/.style={isnode, fill=isOT!15, draw=isOT, thick, minimum width=24mm},
  zone/.style={
    draw, rounded corners=4pt, thick, dashed, inner sep=4mm
  },
  conduit/.style={-{Stealth[length=2.5mm]}, thick, isNet!80!black},
  attack/.style={-{Stealth[length=2.5mm]}, thick, isAttack, dashed},
  flow/.style={-{Stealth[length=2.5mm]}, thick, isOT!70!black},
  netline/.style={thick, black!60},
  axisline/.style={-{Stealth[length=2mm]}, thick, black!70},
  tag/.style={font=\scriptsize\itshape, fill=white, inner sep=1pt},
  level/.style={
    draw, fill=isLight, minimum width=0.85\linewidth, minimum height=8mm,
    font=\small, anchor=west
  },
  brace/.style={decorate, decoration={brace, amplitude=4pt, raise=2pt}},
  legendbox/.style={draw, rounded corners=2pt, fill=isLight, inner sep=2mm, font=\scriptsize, align=left}
}

% --- handy macros ----------------------------------------------------
% \plcicon, \hmiicon etc as simple labelled boxes; useful inline.
\newcommand{\plcicon}[1][PLC]{\tikz[baseline=-0.6ex]{\node[plc, minimum width=10mm, minimum height=5mm, font=\scriptsize, inner sep=1pt]{#1};}}
\newcommand{\hmiicon}[1][HMI]{\tikz[baseline=-0.6ex]{\node[hmi, minimum width=10mm, minimum height=5mm, font=\scriptsize, inner sep=1pt]{#1};}}
\newcommand{\fwicon}[1][FW]{\tikz[baseline=-0.6ex]{\node[firewall, minimum width=10mm, minimum height=5mm, font=\scriptsize, inner sep=1pt]{#1};}}


\providecommand{\courseAuthor}{Industrial Security Lecture Series}
\providecommand{\courseInstitute}{Department of Computer Science}
\providecommand{\companyLogo}{logo}

\definecolor{slideAccent}{HTML}{00205B}
\definecolor{slideSecondary}{HTML}{00A0DC}

% --- Corporate branding overrides ------------------------------------
\IfFileExists{../../../branding.tex}{% =====================================================================
% Corporate / institutional branding -- single file to customise the
% look of the entire course (slides, notes, exercises, labs).
%
% HOW TO REBRAND IN UNDER A MINUTE:
%   1. Replace `branding/logo.pdf` with your own logo
%      (PDF preferred; PNG and JPG also work -- name it `logo.<ext>`).
%   2. (Optional) change the two colours below to your brand palette.
%   3. (Optional) override the author/institute lines below.
%   4. Re-run `make` at the repo root. That's it.
%
% Everything in this file has sensible defaults; you can delete a line
% to fall back to the built-in defaults, or comment it out.
%
% This file is loaded automatically by the slides, exercise, and lab
% preambles -- you never need to \input it manually.
% =====================================================================

% --- Logo --------------------------------------------------------------
% Path is resolved from each leaf-build directory; the default points at
% the shared `branding/` folder at the repo root. If you put your logo
% somewhere else, set the full or relative path here (without extension):
\renewcommand{\companyLogo}{../../../branding/logo}

% --- Author / institute (appears on the title page footer) ------------
\renewcommand{\courseAuthor}{}
\renewcommand{\courseInstitute}{}

% --- Brand colours ----------------------------------------------------
% slideAccent      = primary brand colour (title bands, accent rule)
% slideSecondary   = secondary brand colour (section badge, highlight)
% Both use HTML hex codes. Keep enough contrast against white text.
\definecolor{slideAccent}{HTML}{00205B}      % deep navy
\definecolor{slideSecondary}{HTML}{00A0DC}   % light blue accent
}{}

\colorlet{labAccent}{slideAccent}
\colorlet{labSec}{slideSecondary}
\definecolor{labCheck}{HTML}{2E7D32}
\definecolor{labWarn}{HTML}{B23A48}

\lstset{
  backgroundcolor=\color{black!4},
  basicstyle=\ttfamily\footnotesize,
  breaklines=true,
  frame=leftline,
  framerule=2pt,
  rulecolor=\color{labAccent},
  numbers=left,
  numberstyle=\tiny\color{black!50},
  showstringspaces=false,
  tabsize=2
}

\setlength{\tabcolsep}{4pt}

% Let TeX add a little extra inter-word stretch as a last resort before
% it reports an Overfull \hbox. Only kicks in on lines that would
% otherwise overflow (long \texttt paths, URLs, CIDRs); never loosens a
% line that already fits. Keeps prose/itemize within the margin.
\setlength{\emergencystretch}{3em}

\pagestyle{fancy}
\fancyhf{}
\renewcommand{\headrulewidth}{0.4pt}
\lhead{\small \courseTitle{} -- Lab \labNumber}
\rhead{\small Maps to Section \labSection}
\cfoot{\thepage}

\newtcolorbox{stepbox}[1]{
  enhanced, breakable, arc=2pt, boxrule=0pt,
  colback=labAccent!7, colframe=labAccent, leftrule=3pt,
  fonttitle=\bfseries\small, coltitle=white,
  colbacktitle=labAccent,
  title={\faTools\ Step #1}
}

\newtcolorbox{warnbox}{
  enhanced, breakable, arc=2pt, boxrule=0pt,
  colback=labWarn!7, colframe=labWarn, leftrule=4pt,
  fonttitle=\bfseries\small, coltitle=white, colbacktitle=labWarn,
  title={\faExclamationTriangle\ Caution}
}

\newtcolorbox{notebox}{
  enhanced, breakable, arc=2pt, boxrule=0pt,
  colback=labAccent!4, colframe=labAccent, leftrule=2pt,
  fonttitle=\bfseries\small, coltitle=white, colbacktitle=labAccent,
  title={\faInfoCircle\ Note}
}

\newtcolorbox{goalbox}{
  enhanced, breakable, arc=2pt, boxrule=0pt,
  colback=labCheck!7, colframe=labCheck, leftrule=3pt,
  fonttitle=\bfseries\small, coltitle=white, colbacktitle=labCheck,
  title={\faBullseye\ Goal}
}

\newtcolorbox{deliverbox}{
  enhanced, breakable, arc=2pt, boxrule=0pt,
  colback=labSec!10, colframe=labSec, leftrule=3pt,
  fonttitle=\bfseries\small, coltitle=white, colbacktitle=labSec,
  title={\faClipboardList\ Deliverable}
}

% --- Solution-sheet boxes (used only by lab-solutions.tex) ----------
\newtcolorbox{solutionbox}[1]{
  enhanced, breakable, arc=2pt, boxrule=0pt,
  colback=labCheck!7, colframe=labCheck, leftrule=3pt,
  fonttitle=\bfseries\small, coltitle=white, colbacktitle=labCheck,
  title={\faCheckCircle\ #1}
}

\newtcolorbox{rubricbox}{
  enhanced, breakable, arc=2pt, boxrule=0pt,
  colback=labSec!7, colframe=labSec, leftrule=3pt,
  fonttitle=\bfseries\small, coltitle=white, colbacktitle=labSec,
  title={\faBalanceScale\ Grading rubric}
}

\titleformat{\section}{\large\bfseries\color{labAccent}}{\thesection}{0.5em}{}

\newcommand{\labheader}{%
  \noindent
  \begin{tikzpicture}
    \fill[labAccent] (0,0) rectangle (\textwidth, -1.6cm);
    \node[anchor=north west, text=white, font=\bfseries\large] at (0.6cm, -0.4cm) {Lab \labNumber{} -- \labTitle};
    \node[anchor=south west, text=white, font=\small] at (0.6cm, -1.3cm) {\courseTitle{} \quad $\bullet$ \quad Maps to Section \labSection{} \quad $\bullet$ \quad \labTime};
    \node[anchor=north east, fill=labSec, text=white, font=\bfseries, inner sep=4pt, rounded corners=2pt]
      at ([xshift=-0.4cm, yshift=-0.4cm]\textwidth,0) {LAB \labNumber};
  \end{tikzpicture}
  \vspace{4mm}
}

% =====================================================================

\usepackage[utf8]{inputenc}
\usepackage[T1]{fontenc}
\usepackage[english]{babel}
\usepackage[a4paper, margin=2cm]{geometry}
\usepackage{graphicx}
\usepackage{listings}
\usepackage{xcolor}
\usepackage{colortbl}
\usepackage{tikz}
\usepackage{amsmath}
\usepackage{amssymb}
\usepackage{booktabs}
\usepackage{enumitem}
\usepackage{titlesec}
\usepackage{fancyhdr}
\usepackage{hyperref}
\usepackage{tcolorbox}
\tcbuselibrary{skins, breakable}
\usepackage{fontawesome5}

% =====================================================================
% Shared TikZ styles for the Industrial Security lecture series.
% Loaded by every slide deck and exercise sheet that uses TikZ.
% =====================================================================

\usetikzlibrary{
  arrows.meta,
  shapes,
  shapes.geometric,
  shapes.symbols,
  shapes.multipart,
  positioning,
  fit,
  calc,
  backgrounds,
  decorations.pathmorphing,
  decorations.pathreplacing,
  decorations.text,
  patterns,
  matrix,
  shadows
}

% --- colour palette --------------------------------------------------
\definecolor{isOT}{HTML}{1F4E79}     % OT (deep blue)
\definecolor{isIT}{HTML}{6F2DA8}     % IT (purple)
\definecolor{isSafety}{HTML}{C00000} % safety red
\definecolor{isNet}{HTML}{2E7D32}    % network green
\definecolor{isAttack}{HTML}{B23A48} % attack red
\definecolor{isAccent}{HTML}{E08E0B} % amber
\definecolor{isMute}{HTML}{6B6B6B}   % grey
\definecolor{isLight}{HTML}{EDF2F8}  % light background
\definecolor{isPLC}{HTML}{2C5282}
\definecolor{isHMI}{HTML}{2F855A}
\definecolor{isSCADA}{HTML}{6B46C1}

% --- generic node styles --------------------------------------------
\tikzset{
  isnode/.style={
    draw, rounded corners=2pt, minimum height=8mm, minimum width=18mm,
    font=\small, align=center, thick
  },
  isbox/.style={isnode, fill=isLight},
  plc/.style={isnode, fill=isPLC!15, draw=isPLC, thick},
  hmi/.style={isnode, fill=isHMI!15, draw=isHMI, thick},
  scada/.style={isnode, fill=isSCADA!15, draw=isSCADA, thick},
  sensor/.style={isnode, fill=isNet!10, draw=isNet, minimum height=6mm, minimum width=14mm, font=\scriptsize},
  actuator/.style={isnode, fill=isAccent!15, draw=isAccent, minimum height=6mm, minimum width=14mm, font=\scriptsize},
  firewall/.style={isnode, fill=isAttack!15, draw=isAttack, thick},
  server/.style={isnode, fill=isMute!15, draw=isMute!80, thick},
  switch/.style={isnode, fill=isLight, draw=black!60, minimum height=6mm, minimum width=14mm, font=\scriptsize},
  workstation/.style={isnode, fill=white, draw=isOT, thick},
  attacker/.style={isnode, fill=isAttack!20, draw=isAttack, thick, font=\small\bfseries},
  inetnode/.style={draw, ellipse, minimum width=24mm, minimum height=10mm, font=\scriptsize, fill=isLight, thick},
  process/.style={isnode, fill=isOT!15, draw=isOT, thick, minimum width=24mm},
  zone/.style={
    draw, rounded corners=4pt, thick, dashed, inner sep=4mm
  },
  conduit/.style={-{Stealth[length=2.5mm]}, thick, isNet!80!black},
  attack/.style={-{Stealth[length=2.5mm]}, thick, isAttack, dashed},
  flow/.style={-{Stealth[length=2.5mm]}, thick, isOT!70!black},
  netline/.style={thick, black!60},
  axisline/.style={-{Stealth[length=2mm]}, thick, black!70},
  tag/.style={font=\scriptsize\itshape, fill=white, inner sep=1pt},
  level/.style={
    draw, fill=isLight, minimum width=0.85\linewidth, minimum height=8mm,
    font=\small, anchor=west
  },
  brace/.style={decorate, decoration={brace, amplitude=4pt, raise=2pt}},
  legendbox/.style={draw, rounded corners=2pt, fill=isLight, inner sep=2mm, font=\scriptsize, align=left}
}

% --- handy macros ----------------------------------------------------
% \plcicon, \hmiicon etc as simple labelled boxes; useful inline.
\newcommand{\plcicon}[1][PLC]{\tikz[baseline=-0.6ex]{\node[plc, minimum width=10mm, minimum height=5mm, font=\scriptsize, inner sep=1pt]{#1};}}
\newcommand{\hmiicon}[1][HMI]{\tikz[baseline=-0.6ex]{\node[hmi, minimum width=10mm, minimum height=5mm, font=\scriptsize, inner sep=1pt]{#1};}}
\newcommand{\fwicon}[1][FW]{\tikz[baseline=-0.6ex]{\node[firewall, minimum width=10mm, minimum height=5mm, font=\scriptsize, inner sep=1pt]{#1};}}


\providecommand{\courseAuthor}{Industrial Security Lecture Series}
\providecommand{\courseInstitute}{Department of Computer Science}
\providecommand{\companyLogo}{logo}

\definecolor{slideAccent}{HTML}{00205B}
\definecolor{slideSecondary}{HTML}{00A0DC}

% --- Corporate branding overrides ------------------------------------
\IfFileExists{../../../branding.tex}{% =====================================================================
% Corporate / institutional branding -- single file to customise the
% look of the entire course (slides, notes, exercises, labs).
%
% HOW TO REBRAND IN UNDER A MINUTE:
%   1. Replace `branding/logo.pdf` with your own logo
%      (PDF preferred; PNG and JPG also work -- name it `logo.<ext>`).
%   2. (Optional) change the two colours below to your brand palette.
%   3. (Optional) override the author/institute lines below.
%   4. Re-run `make` at the repo root. That's it.
%
% Everything in this file has sensible defaults; you can delete a line
% to fall back to the built-in defaults, or comment it out.
%
% This file is loaded automatically by the slides, exercise, and lab
% preambles -- you never need to \input it manually.
% =====================================================================

% --- Logo --------------------------------------------------------------
% Path is resolved from each leaf-build directory; the default points at
% the shared `branding/` folder at the repo root. If you put your logo
% somewhere else, set the full or relative path here (without extension):
\renewcommand{\companyLogo}{../../../branding/logo}

% --- Author / institute (appears on the title page footer) ------------
\renewcommand{\courseAuthor}{}
\renewcommand{\courseInstitute}{}

% --- Brand colours ----------------------------------------------------
% slideAccent      = primary brand colour (title bands, accent rule)
% slideSecondary   = secondary brand colour (section badge, highlight)
% Both use HTML hex codes. Keep enough contrast against white text.
\definecolor{slideAccent}{HTML}{00205B}      % deep navy
\definecolor{slideSecondary}{HTML}{00A0DC}   % light blue accent
}{}

\colorlet{labAccent}{slideAccent}
\colorlet{labSec}{slideSecondary}
\definecolor{labCheck}{HTML}{2E7D32}
\definecolor{labWarn}{HTML}{B23A48}

\lstset{
  backgroundcolor=\color{black!4},
  basicstyle=\ttfamily\footnotesize,
  breaklines=true,
  frame=leftline,
  framerule=2pt,
  rulecolor=\color{labAccent},
  numbers=left,
  numberstyle=\tiny\color{black!50},
  showstringspaces=false,
  tabsize=2
}

\setlength{\tabcolsep}{4pt}

% Let TeX add a little extra inter-word stretch as a last resort before
% it reports an Overfull \hbox. Only kicks in on lines that would
% otherwise overflow (long \texttt paths, URLs, CIDRs); never loosens a
% line that already fits. Keeps prose/itemize within the margin.
\setlength{\emergencystretch}{3em}

\pagestyle{fancy}
\fancyhf{}
\renewcommand{\headrulewidth}{0.4pt}
\lhead{\small \courseTitle{} -- Lab \labNumber}
\rhead{\small Maps to Section \labSection}
\cfoot{\thepage}

\newtcolorbox{stepbox}[1]{
  enhanced, breakable, arc=2pt, boxrule=0pt,
  colback=labAccent!7, colframe=labAccent, leftrule=3pt,
  fonttitle=\bfseries\small, coltitle=white,
  colbacktitle=labAccent,
  title={\faTools\ Step #1}
}

\newtcolorbox{warnbox}{
  enhanced, breakable, arc=2pt, boxrule=0pt,
  colback=labWarn!7, colframe=labWarn, leftrule=4pt,
  fonttitle=\bfseries\small, coltitle=white, colbacktitle=labWarn,
  title={\faExclamationTriangle\ Caution}
}

\newtcolorbox{notebox}{
  enhanced, breakable, arc=2pt, boxrule=0pt,
  colback=labAccent!4, colframe=labAccent, leftrule=2pt,
  fonttitle=\bfseries\small, coltitle=white, colbacktitle=labAccent,
  title={\faInfoCircle\ Note}
}

\newtcolorbox{goalbox}{
  enhanced, breakable, arc=2pt, boxrule=0pt,
  colback=labCheck!7, colframe=labCheck, leftrule=3pt,
  fonttitle=\bfseries\small, coltitle=white, colbacktitle=labCheck,
  title={\faBullseye\ Goal}
}

\newtcolorbox{deliverbox}{
  enhanced, breakable, arc=2pt, boxrule=0pt,
  colback=labSec!10, colframe=labSec, leftrule=3pt,
  fonttitle=\bfseries\small, coltitle=white, colbacktitle=labSec,
  title={\faClipboardList\ Deliverable}
}

% --- Solution-sheet boxes (used only by lab-solutions.tex) ----------
\newtcolorbox{solutionbox}[1]{
  enhanced, breakable, arc=2pt, boxrule=0pt,
  colback=labCheck!7, colframe=labCheck, leftrule=3pt,
  fonttitle=\bfseries\small, coltitle=white, colbacktitle=labCheck,
  title={\faCheckCircle\ #1}
}

\newtcolorbox{rubricbox}{
  enhanced, breakable, arc=2pt, boxrule=0pt,
  colback=labSec!7, colframe=labSec, leftrule=3pt,
  fonttitle=\bfseries\small, coltitle=white, colbacktitle=labSec,
  title={\faBalanceScale\ Grading rubric}
}

\titleformat{\section}{\large\bfseries\color{labAccent}}{\thesection}{0.5em}{}

\newcommand{\labheader}{%
  \noindent
  \begin{tikzpicture}
    \fill[labAccent] (0,0) rectangle (\textwidth, -1.6cm);
    \node[anchor=north west, text=white, font=\bfseries\large] at (0.6cm, -0.4cm) {Lab \labNumber{} -- \labTitle};
    \node[anchor=south west, text=white, font=\small] at (0.6cm, -1.3cm) {\courseTitle{} \quad $\bullet$ \quad Maps to Section \labSection{} \quad $\bullet$ \quad \labTime};
    \node[anchor=north east, fill=labSec, text=white, font=\bfseries, inner sep=4pt, rounded corners=2pt]
      at ([xshift=-0.4cm, yshift=-0.4cm]\textwidth,0) {LAB \labNumber};
  \end{tikzpicture}
  \vspace{4mm}
}


\begin{document}
\labheader

\begin{goalbox}
Run a complete Cyber-PHA / Cyber-HAZOP on a SIL-rated Safety Instrumented
System (SIS) per ISA TR84.00.09, quantify the safety integrity it claims
($\mathrm{PFD}_{avg}$, SIL band), and then \emph{break} that claim with a
TRITON/TRISIS-class cyber attack -- showing exactly which LOPA independence
assumption collapses. You will then design security controls that protect
the SIS \emph{without degrading the safety function}, reconcile the
IEC~61511 and IEC~62443 lifecycles on one timeline, and assemble a combined
safety+security assurance case (GSN sketch). Every artefact you produce is
the form an HSE auditor and a NIS2 auditor would each accept.
\end{goalbox}

\section*{Why this lab exists}

In the Bachelor risk-management lab you ran a Cyber-PHA on a toy conveyor
where the ``SIS'' was a hard-wired e-stop. Here the safety function lives
\emph{inside} a programmable logic solver -- a Triconex-class SIS -- and that
changes everything. The classical functional-safety math (IEC~61508/61511)
assumes faults are \emph{random} and \emph{independent}. A cyber adversary is
neither random nor independent: a single logic download can defeat redundancy,
diagnostics, and the protection layers you counted as separate. This lab is
where you learn to do the safety math \emph{and} the attack analysis on the
same piece of equipment, and to reconcile the two governing standards that
most plants run as disconnected silos.

\medskip\noindent
\textbf{Master-level expectation.} No IT analogies, no scaffolded tables with
worked first rows. You are expected to drive the Cyber-HAZOP yourself, to
compute $\mathrm{PFD}_{avg}$ from first principles, and to argue the
independence-defeat formally. Worked artefacts in the solution sheet are the
\emph{reference standard}, not a crutch.

\section*{The scenario -- ``Hydrogen Furnace HF-201''}

\begin{notebox}
A steam-methane-reformer feed line delivers hydrogen to a fired heater
(furnace HF-201). The Burner Management System (BMS) is realised as a
triple-modular-redundant (TMR) safety logic solver (Triconex-class) running
the safety application; a separate BPCS (DCS) runs normal combustion control.
Two safety functions are in scope:

\begin{itemize}[leftmargin=1.4em]\setlength{\itemsep}{1pt}
  \item \textbf{SIF-1 (high pressure):} on $\mathrm{PT}$ high-high in the
        H$_2$ header, close the double-block-and-bleed (DBB) fuel valves
        within the process safety time. Demand mode. Target \textbf{SIL\,2}.
  \item \textbf{SIF-2 (flame failure):} on loss of flame detected by 2oo3
        flame scanners, isolate fuel. Demand mode. Target \textbf{SIL\,3}.
\end{itemize}

Architecture facts you will need:
\begin{itemize}[leftmargin=1.4em]\setlength{\itemsep}{1pt}
  \item SIS engineering workstation (SIS-EWS) runs the vendor TriStation-class
        tool; it is dual-homed (one NIC on the SIS network, one on the plant
        DMZ for the historian feed). The key-switch on the controller has been
        left in PROGRAM for eight months ``for convenience.''
  \item Logic downloads are unsigned; the proprietary engineering protocol has
        no authentication. The SIS shares the safety-network switch with the
        BPCS controllers (flat OT LAN, no zone firewall).
  \item LOPA credits the SIS \emph{and} a BPCS high-pressure alarm + operator
        action \emph{and} a mechanical relief valve as three independent layers.
\end{itemize}
\end{notebox}

\begin{warnbox}
Fictional plant. SIL targets, $\mathrm{PFD}$ figures and component data in
this lab are pedagogical. Do not transplant any number into a real safety
calculation or a real client's safety case -- redact and discard if your
worksheet ever resembles a real asset.
\end{warnbox}

\section*{Part A -- Cyber-HAZOP on the SIS (ISA TR84.00.09)}

\begin{stepbox}{1 -- Node decomposition and the security envelope}
Decompose HF-201 into HAZOP nodes (H$_2$ header up to the DBB valves; the
burner/firebox; the SIS logic solver + I/O; the SIS-EWS and its network
path). For \emph{each} node record the conventional process intent
\emph{and} its security envelope: who can reach it, over which protocol, with
what authentication. Produce a one-paragraph node register. The SIS logic
solver and the SIS-EWS path are the two nodes you will spend the rest of the
lab on -- TR84.00.09 calls these the ``cyber-relevant'' nodes.
\end{stepbox}

\begin{stepbox}{2 -- Guide-word sweep with paired causes}
For the H$_2$ header node and the SIS logic-solver node, walk the IEC~61882
guide words (No, More, Less, Reverse, As~well~as, Part~of, Other~than).
For each \emph{credible} deviation record, in a single worksheet, a row with:
guide word, parameter, deviation, \textbf{one random/systematic safety cause}
\emph{and} \textbf{one cyber cause}, the consequence, the existing safeguard,
and the residual gap. Do not write ``add MFA'' -- name the concrete mechanism
(signed logic download, key-switch interlock, write allow-list on the
engineering protocol, 2oo3 sensor voting integrity check). Aim for at least
six fully-worked rows, of which at least three have a cyber cause that targets
the SIS itself rather than the BPCS. This worksheet is your first deliverable.
\end{stepbox}

\begin{stepbox}{3 -- Classify each cyber cause against the SIS lifecycle phase}
ISA TR84.00.09 maps cyber risk onto the IEC~61511 safety lifecycle phase
where it bites. For every cyber cause in your Step-2 worksheet, tag the phase
it exploits: \emph{design} (e.g.\ flat network, no zones), \emph{realisation}
(unsigned download path), \emph{operation} (key-switch in PROGRAM, shared
logins), or \emph{modification} (uncontrolled MOC bypassing re-validation).
A cause that maps to \emph{operation} or \emph{modification} is the most
dangerous class: it defeats the safety function \emph{after} certification,
so the SIL certificate on the wall is no longer evidence of anything.
\end{stepbox}

\section*{Part B -- The SIL claim and how a compromise voids it}

\begin{stepbox}{4 -- Compute $\mathrm{PFD}_{avg}$ and verify the SIL band}
Take SIF-1 (target SIL\,2) and verify its claimed integrity from component
data. Use the simplified low-demand formula for each subsystem and sum:
\[
\mathrm{PFD}_{avg}^{SIF}
   \;=\; \mathrm{PFD}_S \;+\; \mathrm{PFD}_L \;+\; \mathrm{PFD}_{FE}
\]
where $S$ = sensor subsystem, $L$ = logic solver, $FE$ = final element.
For a simplex (1oo1) element use $\mathrm{PFD}_{avg}\approx
\tfrac{1}{2}\,\lambda_{DU}\,T_I$ and for a 1oo2 redundant element use
$\mathrm{PFD}_{avg}\approx \tfrac{1}{3}\,(\lambda_{DU}\,T_I)^2 +
\beta\,\lambda_{DU}\,\tfrac{T_I}{2}$ (the $\beta$ term is the common-cause
contribution). Component data to use:
\begin{itemize}[leftmargin=1.4em]\setlength{\itemsep}{1pt}
  \item Pressure transmitters: 1oo2, $\lambda_{DU}=300\ \mathrm{FIT}$ each,
        $\beta=10\%$.
  \item TMR logic solver: vendor-certified $\mathrm{PFD}_{avg}=
        2.0\times10^{-4}$.
  \item Final element (DBB valve + solenoid): 1oo1,
        $\lambda_{DU}=600\ \mathrm{FIT}$.
  \item Proof-test interval $T_I = 1\ \mathrm{year} = 8760\ \mathrm{h}$;
        $1\ \mathrm{FIT}=10^{-9}\ \mathrm{h}^{-1}$.
\end{itemize}
Compute each subsystem term, sum, and state the SIL band achieved (SIL\,2 is
$10^{-3}\le \mathrm{PFD}_{avg} < 10^{-2}$, i.e.\ RRF $100$--$1000$). Record
your numbers; the solution sheet gives the reference calculation.
\end{stepbox}

\begin{stepbox}{5 -- Show how a cyber compromise defeats LOPA independence}
LOPA grants risk-reduction credit only to layers that are \emph{independent}
of each other and of the initiating event. Your scenario credits three:
the SIS (SIF-1), the BPCS high-pressure alarm + operator, and the mechanical
relief valve. Now model a TRISIS-class adversary who has reached the SIS-EWS:

\begin{enumerate}[leftmargin=1.6em]\setlength{\itemsep}{1pt}
  \item State, for each of the three layers, whether the adversary can defeat
        it from the SIS-EWS foothold given the architecture facts (flat LAN,
        unsigned download, key-switch in PROGRAM, shared switch with BPCS).
  \item Identify which two layers are \emph{no longer independent} once the
        SIS-EWS is compromised, and explain the mechanism (hint: the same
        engineering path and flat network reach both the SIS logic and the
        BPCS, and the relief valve is the only layer outside the cyber blast
        radius).
  \item Recompute the credited risk reduction \emph{after} the compromise:
        a layer the attacker controls gets a PFD of $1$ (no credit). State the
        residual frequency before and after, and the order-of-magnitude loss
        of protection. This collapse -- not the malware payload -- is the real
        finding. Record the analysis as a short structured argument.
\end{enumerate}
\end{stepbox}

\begin{stepbox}{6 -- TRITON / TRISIS attack walk-through and safe-state analysis}
Walk the documented TRITON/TRISIS kill chain against \emph{this} SIS and map
each stage to a defence you control:
\begin{itemize}[leftmargin=1.4em]\setlength{\itemsep}{1pt}
  \item \textbf{Foothold $\rightarrow$ SIS-EWS:} how the dual-homed EWS bridges
        DMZ to safety network. Which control would have stopped it?
  \item \textbf{EWS $\rightarrow$ controller:} the proprietary protocol +
        key-switch in PROGRAM let logic be appended. What detects this?
  \item \textbf{Payload:} TRITON appended a backdoor to the safety controller's
        firmware/logic, intending to let unsafe states pass while the SIS
        \emph{appeared} healthy. The actual incident tripped the plant because
        a memory-validation mismatch put the TMR controller into a safe fault
        state. Discuss: was the trip ``luck'' or ``design''? What does
        fail-safe TMR behaviour buy you, and what does it \emph{not} buy you
        (a degraded-but-running attack)?
\end{itemize}
State, for each stage, the safe-state implication: does the SIS still reach
its safe state on a real demand after this payload? Justify. This is the
analysis a regulator will ask for after any SIS-touching incident.
\end{stepbox}

\section*{Part C -- Controls that protect safety without degrading it}

\begin{stepbox}{7 -- Design non-interfering security controls}
The cardinal rule of co-engineering: \textbf{a security control must never
degrade, delay, or inhibit the safety function.} For each gap from your
Step-2 worksheet, propose a control and prove it is non-interfering. For each
control state: (a) the threat it removes; (b) why it cannot increase
$\mathrm{PFD}_{avg}$, add latency inside the process safety time, or create a
new spurious-trip path; (c) whether it requires re-validation of the SIF
under IEC~61511. Cover at minimum: key-switch returned to RUN under dual
control; SIS network zoned off the BPCS per IEC~62443-3-2; removal of the
dual-homed NIC (one-way historian feed via data diode); signed/whitelisted
logic downloads; independent code-integrity monitoring of the safety
application. Flag any candidate control you must \emph{reject} because it
would touch the safety path -- e.g.\ an inline IPS that could add latency or
drop a safety I/O packet, or host AV on the logic solver. The rejected list
is as important as the accepted list.
\end{stepbox}

\begin{stepbox}{8 -- Reconcile the IEC 61511 and IEC 62443 lifecycles}
Place the IEC~61511 functional-safety lifecycle and the IEC~62443
security lifecycle on one timeline and identify the shared decision gates.
Produce a mapping table: for each safety phase (HAZID/HAZOP $\rightarrow$
SIF allocation $\rightarrow$ SRS $\rightarrow$ design/realisation $\rightarrow$
V\&V $\rightarrow$ operation/MOC) name the matching IEC~62443 activity
(zone/conduit + risk assessment per 62443-3-2 $\rightarrow$ SL-T allocation
$\rightarrow$ security requirements $\rightarrow$ secure design/3-3 SRs
$\rightarrow$ security V\&V $\rightarrow$ patch/MOC). State explicitly where a
\emph{shared change board} must gate both: any modification to the safety
application is simultaneously a safety MOC \emph{and} a security change, and
neither standard's change process may proceed without the other's sign-off.
\end{stepbox}

\section*{Part D -- (Stretch) The combined assurance case}

\begin{stepbox}{9 -- (Stretch) Build a GSN safety+security case fragment}
Sketch a Goal Structuring Notation (GSN) assurance-case fragment whose top
goal is \emph{``HF-201 SIF-1 achieves SIL\,2 and is not defeatable by a
credible cyber adversary.''} Decompose into the two sub-goals (safety
integrity demonstrated by $\mathrm{PFD}_{avg}$; integrity-of-the-safety-
function preserved against the threat model), the strategy that links them,
the context (threat model, proof-test regime), and the solutions/evidence
(the Step-4 calculation, the Step-7 controls, the Step-8 lifecycle gates).
The point of GSN here is to make the \emph{security argument load-bearing
inside the safety case} -- a SIL certificate that ignores cyber is an
incomplete case. Draw it as an indented goal tree (G/S/C/Sn nodes); a TikZ
diagram is welcome but not required.
\end{stepbox}

\section*{Troubleshooting -- common safety/security conflicts}

\begin{notebox}
\begin{itemize}[leftmargin=1.4em]\setlength{\itemsep}{2pt}
  \item \textbf{``Patch the logic solver now.''} Patching a certified SIS can
        invalidate the safety certificate and trigger full re-validation.
        Resolve via the safety MOC process and compensating controls, not an
        emergency patch window.
  \item \textbf{Inline security monitoring on the safety network.} An IPS or
        DPI box that can drop or delay a packet sits \emph{inside} the process
        safety time budget. Use passive taps / SPAN-port monitoring only;
        never inline on the SIS conduit.
  \item \textbf{Key-switch convenience vs.\ security.} Operations leaves the
        key in PROGRAM to avoid trips during tuning. The fix is procedural
        (dual-control key custody + alarm on PROGRAM state), not a control
        that blocks a genuine emergency download.
  \item \textbf{Spurious trips from a security control.} Any control that can
        force a safe-state trip degrades availability and erodes operator
        trust in the SIS. A non-interfering control must be unable to command
        the safety output.
  \item \textbf{Two registers, two auditors.} Maintaining separate safety and
        security risk registers guarantees a gap falls between them. One
        combined process-hazard-and-cyber register is the only defensible
        artefact.
\end{itemize}
\end{notebox}

\begin{deliverbox}
Hand in a single combined \textbf{safety--security case} pack:
\begin{itemize}[leftmargin=1.4em]\setlength{\itemsep}{1pt}
  \item the Cyber-HAZOP worksheet ($\geq 6$ rows, paired safety+cyber causes,
        each cyber cause tagged to its lifecycle phase) -- Steps 2--3;
  \item the SIF-1 $\mathrm{PFD}_{avg}$ calculation with the SIL band achieved
        -- Step 4;
  \item the LOPA-independence-defeat analysis: which layers stop being
        independent, and the before/after residual frequency -- Step 5;
  \item the TRITON/TRISIS walk-through with the per-stage safe-state
        implication -- Step 6;
  \item the non-interfering controls table (accepted \emph{and} rejected, with
        the re-validation flag) -- Step 7;
  \item the IEC~61511~$\leftrightarrow$~IEC~62443 lifecycle mapping with the
        shared change-board gate -- Step 8;
  \item (stretch) the GSN assurance-case fragment -- Step 9.
\end{itemize}
Plus a one-page chairman's summary naming the top three actionable findings
and, for each, the owner (safety / security / operations) and the next plant
turnaround it is tied to.
\end{deliverbox}

\section*{Further reading}

\begin{notebox}
\begin{itemize}[leftmargin=1.4em]\setlength{\itemsep}{1pt}
  \item IEC~61508 (ed.\,2.0, 2010). \emph{Functional safety of E/E/PE
        safety-related systems} -- the $\mathrm{PFD}_{avg}$ / SIL framework.
  \item IEC~61511 (ed.\,2.0, 2016/17). \emph{Functional safety -- SIS for the
        process industry sector} -- the safety lifecycle and MOC.
  \item ISA~TR84.00.09. \emph{Cybersecurity Related to the Functional Safety
        Lifecycle} -- the document this lab operationalises.
  \item IEC~62443-3-2:2020. \emph{Security risk assessment for system design}
        -- zones, conduits, SL-T allocation.
  \item Dragos. \emph{TRISIS Malware Analysis} (2017) and Mandiant/FireEye.
        \emph{TRITON / Attackers Deploy New ICS Attack Framework} (2017).
  \item GSN Community Standard v3 (2021), Safety-Critical Systems Club --
        for the assurance-case notation.
\end{itemize}
\end{notebox}

\begin{warnbox}
Security must never compromise safety. If, anywhere in this lab, a control you
propose could delay, inhibit, or spuriously trip the safety function, it is
wrong by definition -- redesign it. The safe state always wins; your job is to
make sure a cyber adversary cannot \emph{prevent} the plant from reaching it.
\end{warnbox}

\end{document}
